% Part: first-order-logic
% Chapter: syntax-and-semantics
% Section: first-order-languages

\documentclass[../../../include/open-logic-section]{subfiles}

\begin{document}

\olfileid{fol}{syn}{fol}

\olsection{Basa Logika Tataran Kapisan}


Ungkapan logika tataran kapisan kawangun saka kosakata dhasar kang ngemot
\emph{!!{variable}s}, \emph{!!{constant}s}, \emph{!!{predicate}s}, lan
kadhangkala \emph{!!{function}s}. Saka simbol-simbol kasebut, bebarengan
karo konektif logis, kuantor, lan tandha wacan kayata kurung lan koma,
kawangun \emph{term} lan \emph{!!{formula}s}.

\begin{explain}
Sacara ora resmi, !!{predicate}s iku jeneng kanggo sipat lan relasi,
!!{constant}s iku jeneng kanggo objek siji-siji, lan !!{function}s iku
jeneng kanggo pemetaan. Kajaba !!{identity}~$\eq$, kabeh iki kalebu
\emph{simbol nonlogis} lan bebarengan mbentuk sawijining basa. Saben
basa logika tataran kapisan~$\Lang L$ ditemtokake dening simbol-simbol
nonlogise. Ing kasus kang paling umum, $\Lang L$ ngemot tanpa wates
simbol saka saben jinis.
\end{explain}

Ing kasus umum, kita nggunakake simbol-simbol iki ing logika tataran
kapisan:

\begin{enumerate}
\item Simbol logis
\begin{enumerate}
\item Konektif logis:
  \startycommalist
  \iftag{prvNot}{\ycomma $\lnot$ (negasi)}{}%
  \iftag{prvAnd}{\ycomma $\land$ (konjungsi)}{}%
  \iftag{prvOr}{\ycomma $\lor$ (disjungsi)}{}%
  \iftag{prvIf}{\ycomma $\lif$ (!!{conditional})}{}%
  \iftag{prvIff}{\ycomma $\liff$ (!!{biconditional})}{}%
  \iftag{prvAll}{\ycomma $\lforall$ (kuantor universal)}{}%
  \iftag{prvEx}{\ycomma $\lexists$ (kuantor eksistensial)}{}.
\tagitem{prvFalse}{Konstanta proposisional kanggo !!{falsity}~$\lfalse$.}{}
\tagitem{prvTrue}{Konstanta proposisional kanggo !!{truth}~$\ltrue$.}{}
\item !!{identity} rong panggonan~$\eq$.
\item Himpunan kang !!{denumerable}s saka !!{variable}s: $\Obj v_0$, $\Obj v_1$, $\Obj
  v_2$, \dots
\end{enumerate}
\item Simbol nonlogis kang mbentuk \emph{basa standar} logika tataran
  kapisan
\begin{enumerate}
\item Himpunan kang !!{denumerable}s saka !!{predicate}s $n$-panggonan
  kanggo saben $n>0$: $\Obj
  A^n_0$, $\Obj A^n_1$, $\Obj A^n_2$, \dots
\item Himpunan kang !!{denumerable}s saka !!{constant}s: $\Obj c_0$, $\Obj c_1$, $\Obj
  c_2$, \dots.
\item Himpunan kang !!{denumerable}s saka !!{function}s $n$-panggonan
  kanggo saben $n>0$:
  $\Obj f^n_0$, $\Obj f^n_1$, $\Obj f^n_2$, \dots
\end{enumerate}
\item Tandha wacan: (, ), lan koma.
\end{enumerate}

Umume definisi lan asil kita bakal dirumusake kanggo basa standar
logika tataran kapisan kang jangkep. Nanging, miturut panganggone, kita uga
bisa matesi basa kasebut supaya mung ngemot sawetara !!{predicate}s,
!!{constant}s, lan !!{function}s.

\begin{ex}
Basa aritmetika~$\Lang L_A$ ngemot siji !!{predicate} rong panggonan~$<$,
siji !!{constant}~$\Obj 0$, siji !!{function} siji panggonan~$\prime$,
lan loro !!{function}s rong panggonan~$+$ lan~$\times$.
\end{ex}

\begin{ex}
Basa teori himpunan~$\Lang L_Z$ mung ngemot siji !!{predicate} rong
panggonan~$\in$.
\end{ex}

\begin{ex}
Basa tatanan~$\Lang L_\le$ mung ngemot siji !!{predicate} rong
panggonan~$\le$.
\end{ex}

Maneh, iki mung konvensi: sacara resmi, kabeh iku mung jeneng liya.
Tuladhane, $<$, $\in$, lan $\le$ minangka jeneng liya kanggo $\Obj A^2_0$,
$\Obj 0$ kanggo $\Obj c_0$, $\prime$ kanggo $\Obj f^1_0$, $+$ kanggo
$\Obj f^2_0$, lan $\times$ kanggo $\Obj f^2_1$.

\iftag{defNot,defOr,defAnd,defIf,defIff,defTrue,defFalse,defEx,defAll}{%
Saliyane konektif primitif lan
\iftag{notprvEx,notprvAll}{kuantor}{kuantor-kuantor} kang ditepungake
ing ndhuwur, kita uga nggunakake simbol-simbol \emph{kang ditetepake}
iki:
\startycommalist
  \iftag{defNot}{\ycomma $\lnot$ (negasi)}{}%
  \iftag{defAnd}{\ycomma $\land$ (konjungsi)}{}%
  \iftag{defOr}{\ycomma $\lor$ (disjungsi)}{}%
  \iftag{defIf}{\ycomma $\lif$ (!!{conditional})}{}%
  \iftag{defIff}{\ycomma $\liff$ (!!{biconditional})}{}%
  \iftag{defAll}{\ycomma $\lforall$ (kuantor universal)}{}%
  \iftag{defEx}{\ycomma $\lexists$ (kuantor eksistensial)}{}%
  \iftag{defFalse}{\ycomma !!{falsity}~$\lfalse$}{}%
  \iftag{defTrue}{\ycomma !!{truth}~$\ltrue$}}{}.

\begin{tagblock}{defNot,defOr,defAnd,defIf,defIff,defTrue,defFalse,defEx,defAll}
\begin{explain}
Simbol kang ditetepake dudu perangan resmi saka basa, nanging
ditepungake minangka cekakan ora resmi: kanthi simbol kasebut, kita
bisa nyekakake formula kang bakal dawa banget yen mung nggunakake
simbol primitif. Iki cetha migunani. Nanging paedah kang luwih gedhe
yaiku bukti dadi luwih cekak. Yen sawijining simbol iku primitif,
simbol kasebut kudu ditangani dhewe ing bukti. Mula, saya akeh simbol
primitif, saya dawa bukti kita.
\end{explain}
\end{tagblock}

% Alternate symbols

\begin{intro}
Mbokmenawa sampeyan wis kulina karo istilah lan simbol kang beda saka
kang digunakake ing ndhuwur. Buku teks logika (lan guru) lumrahe
nggunakake $\sim$, $\neg$, utawa~!\ kanggo ``negasi'', lan $\wedge$,
$\cdot$, utawa $\&$ kanggo ``konjungsi''. Simbol kang umum digunakake
kanggo ``kondisional'' utawa ``implikasi'' yaiku $\rightarrow$,
$\Rightarrow$, lan $\supset$.
\iftag{prvIff,defIff}{Simbol kanggo ``bikondisional'', ``bi-implikasi'',
  utawa ``ekuivalensi (material)'' yaiku $\leftrightarrow$,
  $\Leftrightarrow$, lan $\equiv$.}{}
\iftag{prvFalse,defFalse}{Ing pustaka logika basa Inggris, simbol
  $\lfalse$ duwe maneka jeneng: ``falsity'' (nilai bebener luput), ``falsum''
  (istilah Latin), ``absurdity'' (absurditas), utawa ``bottom'' (ngisor).}{}
\iftag{prvTrue,defTrue}{Ing pustaka logika basa Inggris, simbol
  $\ltrue$ duwe maneka jeneng: ``truth'' (nilai bebener bener), ``verum'' (istilah
  Latin), utawa ``top'' (ndhuwur).}{}

Miturut konvensi, aksara cilik saka wiwitan abjad Latin (umpamane
$a$, $b$, $c$) digunakake kanggo !!{constant}s (kadhangkala diarani
jeneng), lan aksara cilik saka pungkasane (umpamane $x$, $y$, $z$)
kanggo !!{variable}s. Kuantor digandhengake karo !!{variable}s,
umpamane $x$; variasi notasi kanggo kuantor universal kalebu $\forall x$,
$(\forall x)$, $(x)$, $\Pi x$, $\bigwedge_x$, dene kanggo kuantor
eksistensial kalebu $\exists x$, $(\exists
x)$, $(Ex)$, $\Sigma x$, $\bigvee_x$.
\end{intro}

\begin{explain}
Kita bisa nganggep kabeh operator proposisional lan loro kuantor
minangka simbol primitif saka basa. Utawa kita bisa milih simbol
primitif kang luwih sithik lan nganggep !!{operator}s liyane minangka
simbol kang ditetepake. Himpunan operator Boolean kang ``jangkep
miturut fungsi bebener'' kalebu $\{ \lnot, \lor \}$, $\{ \lnot,
\land \}$, lan $\{ \lnot,
\lif\}$---saben himpunan iki bisa digandhengake karo salah siji
kuantor kanggo mbentuk basa logika tataran kapisan kang jangkep daya
andharane.

Mbokmenawa sampeyan uga wis ngerti loro !!{operator}s liyane: garis
Sheffer~$|$ (dijenengi miturut Henry Sheffer), lan panah
Peirce~$\downarrow$, kang uga diarani belati Quine. Yen diwenehi
teges kang lumrah, yaiku ``nand'' lan ``nor'' (miturut urutane),
operator-operator iki saben-saben jangkep miturut fungsi bebener.
\end{explain}

\end{document}
